\chapter{BV Integration And Finite-Dimensional Localization}
\label{ch:tqft-bv-integration-localization}

\ConventionsEuclidean

Topological and cohomological QFTs often reduce protected quantities to
finite-dimensional integrals.  This chapter records the precise finite
dimensional BV and localization identities that such reductions rely on, so
that later infinite-dimensional applications can state their additional
analytic hypotheses.

\section{Odd Symplectic Datum}

Let \(\mathcal F\) be a finite-dimensional graded manifold with an odd
symplectic form \(\omega\) of degree \(-1\).  A function \(S\) of degree zero
defines a Hamiltonian vector field \(Q_S\) by
\[
  \iota_{Q_S}\omega=\delta S.
\]
The associated odd Poisson bracket is
\[
  \{F,G\}=\iota_{Q_F}\iota_{Q_G}\omega .
\]
The classical master equation is
\[
  \{S,S\}=0,
\]
which is equivalent to \(Q_S^2=0\).

\section{Gauge Fixing As Lagrangian Choice}

A BV integral is not an integral over all fields.  It is an integral over a
Lagrangian submanifold
\[
  \mathcal L\subset\mathcal F
\]
chosen as a gauge-fixing cycle.  For a half-density \(\rho\) satisfying the
quantum master equation, the expectation of an observable \(\mathcal O\) is
\[
  \langle\mathcal O\rangle_{\mathcal L}
  =
  \int_{\mathcal L}\mathcal O\,\rho .
\]
If \(\mathcal O\) is BV-closed and the deformation of \(\mathcal L\) has no
boundary contribution, the integral is invariant under smooth changes of
\(\mathcal L\).  The proof is Stokes' theorem applied to the odd symplectic
divergence operator.

\section{BV Laplacian Identity}

In Darboux coordinates \((x^i,\xi_i)\) with \(|\xi_i|=1-|x^i|\), the BV
Laplacian on functions relative to the coordinate Berezinian is
\[
  \Delta_{\mathrm{BV}}
  =
  \sum_i(-1)^{|x^i|}
  \frac{\partial^2}{\partial x^i\partial \xi_i}.
\]
The quantum master equation for a density written as
\[
  \rho=\ee^{S/\hbar}
\]
is
\[
  \frac12\{S,S\}+\hbar\Delta_{\mathrm{BV}}S=0.
\]
For any \(F\),
\[
  \Delta_{\mathrm{BV}}(F\ee^{S/\hbar})
  =
  \left(
    \Delta_{\mathrm{BV}}F
    +
    \hbar^{-1}\{S,F\}
  \right)\ee^{S/\hbar}.
\]
Integrating a BV-exact term over a compact Lagrangian gives zero.  This is
the algebraic identity behind Ward identities in finite-dimensional BV
models.

\section{Cohomological Localization Identity}

Let \(M\) be an oriented finite-dimensional manifold with an odd vector field
or differential \(Q\) acting on forms, and suppose
\[
  Q^2=0
\]
on the integrand.  If \(S\) and \(\mathcal O\) are \(Q\)-closed and \(V\) is
an odd functional, define
\[
  I(t)=\int_M \mathcal O\,\ee^{-S-tQV}.
\]
Then
\[
  \frac{dI}{dt}
  =
  -\int_M Q\!\left(V\mathcal O\,\ee^{-S-tQV}\right).
\]
When the integral of a \(Q\)-exact form vanishes, \(I(t)\) is independent of
\(t\).  If the real part of \(QV\) is nonnegative and its zero locus is a
compact submanifold \(F\), the limit \(t\to\infty\) localizes the integral to
\(F\).

\section{One-Loop Normal Form}

Suppose \(F\subset M\) is a smooth fixed locus and the normal quadratic form
of \(QV\) is nondegenerate.  Locally write normal coordinates \(y\) and
fermionic partners \(\chi\).  The quadratic contribution has the form
\[
  t\,y^TAy+t\,\chi^TB\chi .
\]
The Gaussian normal integral produces
\[
  \frac{\operatorname{Pf}(B)}{\sqrt{\det A}},
\]
with orientation and phase determined by the integration cycle.  Thus the
localized integral is
\[
  I=\int_F
  \frac{\iota_F^*(\mathcal O\,\ee^{-S})}
  {e_Q(N_F)},
\]
where \(e_Q(N_F)\) is the equivariant Euler class of the normal complex.

\section{Boundary And Noncompact Terms}

If \(M\) is noncompact or has boundary, the derivative \(dI/dt\) equals a
boundary contribution.  A localization formula must therefore declare one of:
\[
\begin{gathered}
  \text{compact integration cycle},\qquad
  \text{decay estimate at infinity},\\
  \text{boundary condition cancelling the term},\qquad
  \text{explicit residue prescription}.
\end{gathered}
\]
In gauge-theoretic localization, noncompact Coulomb branches, reducible
connections, small instantons, and singular monopole strata are precisely
places where this boundary analysis enters.

\begin{openproblem}[Infinite-dimensional BV localization]
\label{op:infinite-dimensional-bv-localization}
Give theorem-level localization statements for cohomological gauge theories
from a regulated BV-BFV functional integral, including integration cycles,
compactness or boundary estimates, zero modes, determinant lines, singular
strata, and gluing compatibility.
\end{openproblem}
